\documentclass[12pt,a4paper]{ctexart}
\usepackage[UTF8]{ctex}
\usepackage{geometry}
\usepackage{tcolorbox}
\tcbuselibrary{breakable}
\usepackage{tabularx}
\usepackage{array}
\usepackage{amsmath}
\usepackage{titlesec}
\usepackage{graphicx}
\usepackage{caption}
\usepackage{float}
\usepackage{fancyhdr}
\pagestyle{fancy}
\fancyhf{}                % 清空默认页眉页脚
\renewcommand{\headrulewidth}{0pt} % 去掉页眉横线

% ===== 可修改的变量（在此处设置）=====
\newcommand{\expname}{体验实验验证弱相互作用下宇称不守恒}                 
\newcommand{\studentname}{武大兵}              
\newcommand{\grade}{2077}                      
\newcommand{\studentid}{2077486050814}         
\newcommand{\expdate}{2077年 8 月 24 日}        
% ========================================

\geometry{left=0.5cm,right=0.5cm,top=1cm,bottom=1cm}


\raggedbottom
\setlength{\parindent}{2em}

\ctexset{
  section = {
    afterindent = true
  }
}

\titleformat{\section}
{\zihao{4}\bfseries}
{\thesection}
{0.5em}
{}

\renewcommand{\arraystretch}{1.2}

\newcommand{\expimg}[2]{
\begin{figure}[H]
\centering
\includegraphics[width=0.5\linewidth]{#1}
\caption{#2}
\end{figure}
}

\ctexset{
  section = {
    format = \zihao{4}\bfseries,      % 四号加粗
  },
  subsection = {
    format = \zihao{-4}\bfseries,     % 小四加粗
  },
  subsubsection = {
    format = \zihao{-4}\normalfont,    % 小四常规
  }
}

\begin{document}
\thispagestyle{fancy}

\begin{tcolorbox}[
    colback=white, colframe=black,
    boxrule=0.5pt, sharp corners,
    left=0.5cm, right=0.5cm, top=0.5cm, bottom=0.5cm,
    before skip=0pt, after skip=0pt,
    breakable
]

\vspace*{0pt}

% 标题部分
\begin{center}
    \zihao{2} \textbf{武汉大学物理科学与技术学院}\\[0.2cm]
    \zihao{2} \textbf{物理实验报告}\\[0.2cm]
    \zihao{4} 物理科学与技术学院 \quad 物理学专业 \quad \quad \expdate
\end{center}

% 表格
\begin{tabularx}{\linewidth}{|>{\centering\arraybackslash}p{1.8cm}|>{\centering\arraybackslash}X|>{\centering\arraybackslash}p{1cm}|>{\centering\arraybackslash}X|>{\centering\arraybackslash}p{1cm}|>{\centering\arraybackslash}p{3.5cm}|>{\centering\arraybackslash}p{1cm}|>{\centering\arraybackslash}X|}
\hline
\textbf{实验名称} & \multicolumn{7}{>{\centering\arraybackslash}p{0.88\linewidth}|}{\expname} \\
\hline
\textbf{姓 名} &\studentname& \textbf{年 级} &\grade& \textbf{学 号} &\studentid& \textbf{成 绩} & \\
\hline
\end{tabularx}

\vspace{0.5cm}

% 目录结构
\noindent \textbf{实验报告内容：}\\[0.4cm]
\noindent
\begin{minipage}[t]{0.45\textwidth}
    一、实验目的 \\[0.5cm]
    二、主要实验仪器 \\[0.5cm]
    三、实验原理 \\[0.5cm]
    四、实验内容与步骤
\end{minipage}
\hfill
\begin{minipage}[t]{0.45\textwidth}
    五、数据表格 \\[0.5cm]
    六、数据处理及结果表达 \\[0.5cm]
    七、实验结果分析 \\[0.5cm]
    八、习题
\end{minipage}

\vspace{0.5cm}
\noindent\rule{\textwidth}{0.4pt}

% ===== 正文 =====
\setlength{\parindent}{2em}
\setlength{\parskip}{0pt}

\section{实验目的}
拦路雨偏似雪花 饮泣的你冻吗 这风褛我给你磨到有襟花

连调了职也不怕 怎么始终牵挂 苦心选中今天想车你回家
\section{主要实验仪器}
原谅我不再送花 伤口应要结疤

花瓣铺满心底坟场才害怕

如若你非我不嫁 彼此终必火化

一生一世等一天需要代价

\section{实验原理}
谁都只得那双手 靠拥抱亦难任你拥有

要拥有必先懂失去怎接受

曾沿着雪路漫游 为何为好事泪流

谁能凭爱意任富士山私有

\subsection{爱情转移}
回忆是抓不到的月光握紧就变黑暗

让虚假的背影消失于晴朗

阳光在身上流转 等所有业障被原谅

\subsection{爱情不停站}
\subsubsection{想开往地老天荒}
需要多勇敢

\begin{equation}
    E= mc^2
\end{equation}
\section{实验内容与步骤}
情人节不要说穿 只敢抚你发端

这种姿态可会令你更心酸

留在汽车里取暖 应该怎么规劝

怎么可以将手腕忍痛划损
\expimg{figs/111}{CERN}
\section{数据处理及结果表达}
感情需要人接班 接近换来期望

期望带来失望的恶性循环

短暂的总是浪漫 漫长总会不满

烧完美好青春换一个老伴
\section{实验结果分析}
把一个人的温暖转移到另一个的胸膛

让上次犯的错反省出梦想

每个人都是这样

享受过提心吊胆

才拒绝做爱情待罪的羔羊

\section{习题}
何不把悲哀感觉假设是来自你虚构

试管里找不到它染污眼眸

前尘硬化像石头 随缘地抛下便逃走

\subsection{我绝不罕有}
\subsubsection{往街里绕过一周}
我便化乌有

\vspace{1\baselineskip}

你还嫌不够

我把这陈年风褛

送赠你解咒
\end{tcolorbox}

\end{document}